\documentclass[11pt,a4paper]{article}
\usepackage[utf8]{inputenc}
\usepackage[T1]{fontenc}
\usepackage[spanish,es-minimal,es-noshorthands]{babel}
\AtBeginDocument{\renewcommand{\sin}{\operatorname{sen}}\renewcommand{\arcsin}{\operatorname{arcsen}}}
\usepackage{lmodern}
\usepackage[sfdefault,lining]{FiraSans}
\usepackage[a4paper,top=2.3cm,bottom=2.2cm,left=1.8cm,right=1.8cm,headheight=1.3cm,headsep=0.4cm,footskip=1cm]{geometry}
\usepackage{amsmath,amssymb}
\usepackage{icomma}
\usepackage{graphicx}
\usepackage[table]{xcolor}
\usepackage{tikz}
\usepackage{fontawesome5}
\usepackage{ccicons}
\usepackage{fancyhdr}
\usepackage{lastpage}
\usepackage{tasks}
\usepackage{needspace}
\usepackage[most]{tcolorbox}
\usepackage{array}
\usepackage[hidelinks]{hyperref}
\AtBeginDocument{\hypersetup{pdftitle={Ficha de ejercicios de la unidad 1: Números naturales, enteros y fraccionarios. 4º ESO},pdfsubject={Matemáticas A 4º ESO},pdfkeywords={números naturales, operaciones combinadas, potencias, raíces, números enteros, valor absoluto, fracciones, fracciones equivalentes, simplificación, amplificación, problemas con fracciones, propiedades de las potencias}}}

\definecolor{navy}{HTML}{1B2748}
\definecolor{accent}{HTML}{E8603A}
\definecolor{teal}{HTML}{1F9E90}
\definecolor{soft}{HTML}{EEF1F8}
\definecolor{muted}{HTML}{5D6685}

\setlength{\parindent}{0pt}
\setlength{\parskip}{3pt}
\settasks{label=\alph*), label-format=\bfseries\color{navy}, label-width=1.3em, item-indent=1.8em, column-sep=0.6em, after-item-skip=2pt, before-skip=3pt, after-skip=4pt}

% ---------- capçalera i peu ----------
\newcommand{\llicencia}{\href{https://creativecommons.org/licenses/by-nc-sa/4.0/deed.es}{\ccbyncsa}\ \ Licencia Creative Commons Atribución-NoComercial-CompartirIgual 4.0 Internacional (CC BY-NC-SA 4.0)}
\pagestyle{fancy}
\fancyhf{}
\fancyhead[L]{\raisebox{-0.35\height}{\includegraphics[height=0.95cm]{logo-gva.png}}}
\fancyhead[C]{\small\color{muted}Ficha de ejercicios · Unidad 1: Números naturales, enteros y fraccionarios · 4\textsuperscript{o} ESO A}
\fancyhead[R]{\raisebox{-0.35\height}{\includegraphics[height=0.95cm]{logo-vermellar.png}}}
\fancyfoot[L]{\scriptsize\color{muted}\llicencia}
\fancyfoot[R]{\small\color{muted}\thepage}
\renewcommand{\headrulewidth}{0.4pt}
\renewcommand{\headrule}{\color{muted!40}\hrule height \headrulewidth}
\fancypagestyle{primera}{\fancyhf{}\renewcommand{\headrulewidth}{0pt}
  \fancyfoot[L]{\scriptsize\color{muted}\llicencia}
  \fancyfoot[R]{\small\color{muted}\thepage}}

% ---------- nivells i exercicis ----------
\newcommand{\estrelles}[1]{{\color{accent}\ifcase#1\or\faStar\faStar[regular]\faStar[regular]\or\faStar\faStar\faStar[regular]\or\faStar\faStar\faStar\fi}}
\newcounter{exer}
\newcommand{\ex}[1]{\par\needspace{4\baselineskip}\medskip\stepcounter{exer}\noindent{\color{navy}\bfseries\theexer.}\ \ {\small\estrelles{#1}}\ \ }

\newcommand{\seccio}[2]{\par\needspace{10\baselineskip}\bigskip
  \begin{tcolorbox}[colback=navy,colframe=navy,coltext=white,arc=2mm,boxrule=0pt,left=3mm,right=3mm,top=1.2mm,bottom=1.2mm]
  {\large\bfseries #1\quad #2}
  \end{tcolorbox}\vspace{-2pt}}

\newcolumntype{C}[1]{>{\centering\arraybackslash}p{#1}}

% quadre de recordatori (només a la fitxa d'exercicis)
\newcommand{\recorda}[1]{\par\smallskip\begin{tcolorbox}[colback=accent!7,colframe=accent!55,arc=2mm,boxrule=0.5pt,left=2.5mm,right=2.5mm,top=1mm,bottom=1mm]\small{\color{accent}\faLightbulb[regular]}\ \textbf{Recuerda:}\ #1\end{tcolorbox}}
\newcommand{\res}[1]{#1}

\def\fmneed{11}% que el títol d'una secció no quede sol al final del full
% =====================================================================
% Capçalera FentMates de les fitxes: un color per a cada curs.
% Els colors són els de la pàgina inicial (LEVELS a assets/js/materials.js);
% el logo ve de _eines/logo/pdf/ (python3 _eines/logo/genera.py).
%
% Ús, al final del preàmbul de la unitat:
%   % =====================================================================
% Capçalera FentMates de les fitxes: un color per a cada curs.
% Els colors són els de la pàgina inicial (LEVELS a assets/js/materials.js);
% el logo ve de _eines/logo/pdf/ (python3 _eines/logo/genera.py).
%
% Ús, al final del preàmbul de la unitat:
%   % =====================================================================
% Capçalera FentMates de les fitxes: un color per a cada curs.
% Els colors són els de la pàgina inicial (LEVELS a assets/js/materials.js);
% el logo ve de _eines/logo/pdf/ (python3 _eines/logo/genera.py).
%
% Ús, al final del preàmbul de la unitat:
%   \input{../comu/fentmates.tex}
%   \fmcurs{1eso}
% i, al començament del document:
%   \capcalera{Matemàtiques}{1r ESO}{Fitxa d'exercicis · Unitat 2}{Divisibilitat}
% Opcional abans de \input: \def\fmneed{14} (espai mínim abans de cada secció).
% =====================================================================
\makeatletter
\def\fm@logodir{../../logo/pdf}
\providecommand{\fmneed}{10}
\def\fm@cursnom{}

% clau del curs, color i color del text damunt del color
\newcommand{\fm@curs}[3]{\@namedef{fm@color@#1}{#2}\@namedef{fm@text@#1}{#3}}
\fm@curs{1eso}{E6B84A}{navy}
\fm@curs{2eso}{3B82F6}{white}
\fm@curs{3eso}{FF7A50}{white}
\fm@curs{4eso-a}{A78BFA}{white}
\fm@curs{4eso-b}{6D28D9}{white}
\fm@curs{1batx-mi}{9B1C1C}{white}
\fm@curs{1batx-ccss}{EF6B6B}{white}
\fm@curs{2batx-mii}{1E6B47}{white}
\fm@curs{2batx-ccss}{2BB3A3}{white}

\newcommand{\fmcurs}[1]{%
  \@ifundefined{fm@color@#1}{\PackageError{fentmates}{Curs desconegut: #1}{}}{}%
  \edef\fm@tmp{\noexpand\definecolor{curs}{HTML}{\@nameuse{fm@color@#1}}%
               \noexpand\colorlet{curstext}{\@nameuse{fm@text@#1}}}\fm@tmp
  \colorlet{cursfosc}{curs!72!black}%
  \def\fm@logo{\fm@logodir/vertical-#1.pdf}%
  \def\fm@marca{\fm@logodir/marca-#1.pdf}}

% ---------- capçalera de la primera pàgina ----------
% #1 matèria (petit), #2 curs, #3 tipus de document i unitat, #4 títol
\newcommand{\capcalera}[4]{%
  \gdef\fm@cursnom{#2}%
  \par\noindent
  \begin{tcolorbox}[colback=curs,colframe=curs,boxrule=0pt,arc=2.2mm,width=3.5cm,height=2.15cm,
      valign=center,left=3mm,right=2mm,top=1mm,bottom=1mm,nobeforeafter,box align=center]
    \color{curstext}{\footnotesize #1}\par\vspace{1pt}{\fontsize{19}{21}\selectfont\bfseries #2\par}%
  \end{tcolorbox}\hspace{5mm}%
  \begin{minipage}[c]{\dimexpr\linewidth-3.5cm-5mm-2.9cm\relax}
    {\color{cursfosc}\small\bfseries\MakeUppercase{#3}}\par\vspace{2pt}
    {\color{navy}\fontsize{20}{23}\selectfont\bfseries #4\par}
  \end{minipage}\hfill
  \begin{minipage}[c]{2.6cm}\raggedleft\includegraphics[width=2.6cm]{\fm@logo}\end{minipage}\par
  \vspace{2pt}}

% ---------- barres de secció del color del curs ----------
\renewcommand{\seccio}[2]{\par\needspace{\fmneed\baselineskip}\bigskip
  \begin{tcolorbox}[colback=curs,colframe=curs,coltext=curstext,arc=2mm,boxrule=0pt,left=3mm,right=3mm,top=1.2mm,bottom=1.2mm]
  {\large\bfseries #1\quad #2}
  \end{tcolorbox}\vspace{-2pt}}

% ---------- estrelles de dificultat del color del curs (a Batxillerat no se'n posen) ----------
\@ifundefined{estrelles}{}{%
  \renewcommand{\estrelles}[1]{{\color{curs}\ifcase#1\or\faStar\faStar[regular]\faStar[regular]\or\faStar\faStar\faStar[regular]\or\faStar\faStar\faStar\fi}}}

% ---------- capçalera i peu de les pàgines ----------
% primera pàgina: logos institucionals al peu
\fancypagestyle{primera}{\fancyhf{}\renewcommand{\headrulewidth}{0pt}
  \fancyfoot[L]{\raisebox{-0.32\height}{\includegraphics[height=0.78cm]{logo-gva.png}}\hspace{5mm}%
    \raisebox{-0.32\height}{\includegraphics[height=0.78cm]{logo-vermellar.png}}}
  \fancyfoot[R]{\scriptsize\color{muted}\href{https://creativecommons.org/licenses/by-nc-sa/4.0/deed.ca}{\ccbyncsa}\ \ CC BY-NC-SA 4.0 · fentmates.org\qquad\small\thepage}}
% la resta: el símbol a l'esquerra i el curs a la dreta, amb el títol al centre com abans
\fancyhead[L]{\raisebox{-0.3\height}{\includegraphics[height=0.62cm]{\fm@marca}}}
\fancyhead[R]{\small\bfseries\color{cursfosc}\fm@cursnom}
\renewcommand{\headrule}{\color{curs!45}\hrule height \headrulewidth}
\makeatother

%   \fmcurs{1eso}
% i, al començament del document:
%   \capcalera{Matemàtiques}{1r ESO}{Fitxa d'exercicis · Unitat 2}{Divisibilitat}
% Opcional abans de \input: \def\fmneed{14} (espai mínim abans de cada secció).
% =====================================================================
\makeatletter
\def\fm@logodir{../../logo/pdf}
\providecommand{\fmneed}{10}
\def\fm@cursnom{}

% clau del curs, color i color del text damunt del color
\newcommand{\fm@curs}[3]{\@namedef{fm@color@#1}{#2}\@namedef{fm@text@#1}{#3}}
\fm@curs{1eso}{E6B84A}{navy}
\fm@curs{2eso}{3B82F6}{white}
\fm@curs{3eso}{FF7A50}{white}
\fm@curs{4eso-a}{A78BFA}{white}
\fm@curs{4eso-b}{6D28D9}{white}
\fm@curs{1batx-mi}{9B1C1C}{white}
\fm@curs{1batx-ccss}{EF6B6B}{white}
\fm@curs{2batx-mii}{1E6B47}{white}
\fm@curs{2batx-ccss}{2BB3A3}{white}

\newcommand{\fmcurs}[1]{%
  \@ifundefined{fm@color@#1}{\PackageError{fentmates}{Curs desconegut: #1}{}}{}%
  \edef\fm@tmp{\noexpand\definecolor{curs}{HTML}{\@nameuse{fm@color@#1}}%
               \noexpand\colorlet{curstext}{\@nameuse{fm@text@#1}}}\fm@tmp
  \colorlet{cursfosc}{curs!72!black}%
  \def\fm@logo{\fm@logodir/vertical-#1.pdf}%
  \def\fm@marca{\fm@logodir/marca-#1.pdf}}

% ---------- capçalera de la primera pàgina ----------
% #1 matèria (petit), #2 curs, #3 tipus de document i unitat, #4 títol
\newcommand{\capcalera}[4]{%
  \gdef\fm@cursnom{#2}%
  \par\noindent
  \begin{tcolorbox}[colback=curs,colframe=curs,boxrule=0pt,arc=2.2mm,width=3.5cm,height=2.15cm,
      valign=center,left=3mm,right=2mm,top=1mm,bottom=1mm,nobeforeafter,box align=center]
    \color{curstext}{\footnotesize #1}\par\vspace{1pt}{\fontsize{19}{21}\selectfont\bfseries #2\par}%
  \end{tcolorbox}\hspace{5mm}%
  \begin{minipage}[c]{\dimexpr\linewidth-3.5cm-5mm-2.9cm\relax}
    {\color{cursfosc}\small\bfseries\MakeUppercase{#3}}\par\vspace{2pt}
    {\color{navy}\fontsize{20}{23}\selectfont\bfseries #4\par}
  \end{minipage}\hfill
  \begin{minipage}[c]{2.6cm}\raggedleft\includegraphics[width=2.6cm]{\fm@logo}\end{minipage}\par
  \vspace{2pt}}

% ---------- barres de secció del color del curs ----------
\renewcommand{\seccio}[2]{\par\needspace{\fmneed\baselineskip}\bigskip
  \begin{tcolorbox}[colback=curs,colframe=curs,coltext=curstext,arc=2mm,boxrule=0pt,left=3mm,right=3mm,top=1.2mm,bottom=1.2mm]
  {\large\bfseries #1\quad #2}
  \end{tcolorbox}\vspace{-2pt}}

% ---------- estrelles de dificultat del color del curs (a Batxillerat no se'n posen) ----------
\@ifundefined{estrelles}{}{%
  \renewcommand{\estrelles}[1]{{\color{curs}\ifcase#1\or\faStar\faStar[regular]\faStar[regular]\or\faStar\faStar\faStar[regular]\or\faStar\faStar\faStar\fi}}}

% ---------- capçalera i peu de les pàgines ----------
% primera pàgina: logos institucionals al peu
\fancypagestyle{primera}{\fancyhf{}\renewcommand{\headrulewidth}{0pt}
  \fancyfoot[L]{\raisebox{-0.32\height}{\includegraphics[height=0.78cm]{logo-gva.png}}\hspace{5mm}%
    \raisebox{-0.32\height}{\includegraphics[height=0.78cm]{logo-vermellar.png}}}
  \fancyfoot[R]{\scriptsize\color{muted}\href{https://creativecommons.org/licenses/by-nc-sa/4.0/deed.ca}{\ccbyncsa}\ \ CC BY-NC-SA 4.0 · fentmates.org\qquad\small\thepage}}
% la resta: el símbol a l'esquerra i el curs a la dreta, amb el títol al centre com abans
\fancyhead[L]{\raisebox{-0.3\height}{\includegraphics[height=0.62cm]{\fm@marca}}}
\fancyhead[R]{\small\bfseries\color{cursfosc}\fm@cursnom}
\renewcommand{\headrule}{\color{curs!45}\hrule height \headrulewidth}
\makeatother

%   \fmcurs{1eso}
% i, al començament del document:
%   \capcalera{Matemàtiques}{1r ESO}{Fitxa d'exercicis · Unitat 2}{Divisibilitat}
% Opcional abans de \input: \def\fmneed{14} (espai mínim abans de cada secció).
% =====================================================================
\makeatletter
\def\fm@logodir{../../logo/pdf}
\providecommand{\fmneed}{10}
\def\fm@cursnom{}

% clau del curs, color i color del text damunt del color
\newcommand{\fm@curs}[3]{\@namedef{fm@color@#1}{#2}\@namedef{fm@text@#1}{#3}}
\fm@curs{1eso}{E6B84A}{navy}
\fm@curs{2eso}{3B82F6}{white}
\fm@curs{3eso}{FF7A50}{white}
\fm@curs{4eso-a}{A78BFA}{white}
\fm@curs{4eso-b}{6D28D9}{white}
\fm@curs{1batx-mi}{9B1C1C}{white}
\fm@curs{1batx-ccss}{EF6B6B}{white}
\fm@curs{2batx-mii}{1E6B47}{white}
\fm@curs{2batx-ccss}{2BB3A3}{white}

\newcommand{\fmcurs}[1]{%
  \@ifundefined{fm@color@#1}{\PackageError{fentmates}{Curs desconegut: #1}{}}{}%
  \edef\fm@tmp{\noexpand\definecolor{curs}{HTML}{\@nameuse{fm@color@#1}}%
               \noexpand\colorlet{curstext}{\@nameuse{fm@text@#1}}}\fm@tmp
  \colorlet{cursfosc}{curs!72!black}%
  \def\fm@logo{\fm@logodir/vertical-#1.pdf}%
  \def\fm@marca{\fm@logodir/marca-#1.pdf}}

% ---------- capçalera de la primera pàgina ----------
% #1 matèria (petit), #2 curs, #3 tipus de document i unitat, #4 títol
\newcommand{\capcalera}[4]{%
  \gdef\fm@cursnom{#2}%
  \par\noindent
  \begin{tcolorbox}[colback=curs,colframe=curs,boxrule=0pt,arc=2.2mm,width=3.5cm,height=2.15cm,
      valign=center,left=3mm,right=2mm,top=1mm,bottom=1mm,nobeforeafter,box align=center]
    \color{curstext}{\footnotesize #1}\par\vspace{1pt}{\fontsize{19}{21}\selectfont\bfseries #2\par}%
  \end{tcolorbox}\hspace{5mm}%
  \begin{minipage}[c]{\dimexpr\linewidth-3.5cm-5mm-2.9cm\relax}
    {\color{cursfosc}\small\bfseries\MakeUppercase{#3}}\par\vspace{2pt}
    {\color{navy}\fontsize{20}{23}\selectfont\bfseries #4\par}
  \end{minipage}\hfill
  \begin{minipage}[c]{2.6cm}\raggedleft\includegraphics[width=2.6cm]{\fm@logo}\end{minipage}\par
  \vspace{2pt}}

% ---------- barres de secció del color del curs ----------
\renewcommand{\seccio}[2]{\par\needspace{\fmneed\baselineskip}\bigskip
  \begin{tcolorbox}[colback=curs,colframe=curs,coltext=curstext,arc=2mm,boxrule=0pt,left=3mm,right=3mm,top=1.2mm,bottom=1.2mm]
  {\large\bfseries #1\quad #2}
  \end{tcolorbox}\vspace{-2pt}}

% ---------- estrelles de dificultat del color del curs (a Batxillerat no se'n posen) ----------
\@ifundefined{estrelles}{}{%
  \renewcommand{\estrelles}[1]{{\color{curs}\ifcase#1\or\faStar\faStar[regular]\faStar[regular]\or\faStar\faStar\faStar[regular]\or\faStar\faStar\faStar\fi}}}

% ---------- capçalera i peu de les pàgines ----------
% primera pàgina: logos institucionals al peu
\fancypagestyle{primera}{\fancyhf{}\renewcommand{\headrulewidth}{0pt}
  \fancyfoot[L]{\raisebox{-0.32\height}{\includegraphics[height=0.78cm]{logo-gva.png}}\hspace{5mm}%
    \raisebox{-0.32\height}{\includegraphics[height=0.78cm]{logo-vermellar.png}}}
  \fancyfoot[R]{\scriptsize\color{muted}\href{https://creativecommons.org/licenses/by-nc-sa/4.0/deed.ca}{\ccbyncsa}\ \ CC BY-NC-SA 4.0 · fentmates.org\qquad\small\thepage}}
% la resta: el símbol a l'esquerra i el curs a la dreta, amb el títol al centre com abans
\fancyhead[L]{\raisebox{-0.3\height}{\includegraphics[height=0.62cm]{\fm@marca}}}
\fancyhead[R]{\small\bfseries\color{cursfosc}\fm@cursnom}
\renewcommand{\headrule}{\color{curs!45}\hrule height \headrulewidth}
\makeatother

\fmcurs{4eso-a}

\begin{document}
\thispagestyle{primera}

\capcalera{Matemáticas A}{4\textsuperscript{o} ESO}{Ficha de ejercicios · Unidad 1}{\raggedright Números naturales, enteros y fraccionarios}
\vspace{10pt}
Nombre y apellidos: \rule{6.6cm}{0.4pt}\hfill Grupo: \rule{1.6cm}{0.4pt}\hfill Fecha: \rule{2.2cm}{0.4pt}
\vspace{8pt}

\begin{tcolorbox}[colback=soft,colframe=soft,arc=2mm,left=3mm,right=3mm,top=1.5mm,bottom=1.5mm]
\small Ejercicios graduados por niveles de dificultad:\quad \estrelles{1}\ básico \qquad \estrelles{2}\ medio \qquad \estrelles{3}\ avanzado.\\[2pt] Resuélvelos en el cuaderno mostrando todos los pasos. Haz los cálculos sin calculadora.
\end{tcolorbox}

% =====================================================
\seccio{1}{Números naturales. Operaciones combinadas}

\ex{1} Calcula:\begin{tasks}(6)
  \task $3^{5}$
  \task $10^{6}$
  \task $2^{3}\cdot 5^{2}$
  \task $\sqrt{144}$
  \task $\sqrt{225}$
  \task $17^{0}$
\end{tasks}

\ex{1} Calcula el cociente y el resto de cada división y haz la prueba ($D=d\cdot q+r$):\begin{tasks}(2)
  \task $1250:37$
  \task $5000:123$
\end{tasks}

\ex{2} Calcula paso a paso:\begin{tasks}[after-item-skip=6pt](2)
  \task $25-3\cdot \left[12-2\cdot \left(7-4\right)\right]$
  \task $6+48:\left[20-4\cdot \left(9-6\right)\right]$
  \task $60:\left[3\cdot \left(9-5\right)+2\cdot \left(7-3\right)\right]$
  \task $2\cdot \left(40-3\cdot \left[5+2\cdot \left(9-6\right)\right]\right)+7$
  \task $5\cdot \left[18-\left(4+2\cdot 3\right)\right]-4\cdot \left[10-\left(9-2\cdot 3\right)\right]$
  \task $2^{3}\cdot \left[15-\left(\sqrt{49}+2\right)\right]+36:3^{2}$
\end{tasks}\recorda{Primero los paréntesis y corchetes (de dentro hacia fuera); después las potencias y las raíces; después las multiplicaciones y divisiones (de izquierda a derecha) y, al final, las sumas y restas.}

\ex{2} Compramos 3 camisetas a 18\,€ y 2 pantalones a 35\,€, y nos hacen un descuento de 15\,€ sobre el total. Pagamos con dos billetes de 100\,€. Escribe la operación en una sola expresión y calcula el cambio.

\ex{3} Un alumno ha hecho $36:6\cdot2+4^2=36:12+16=3+16=19$. Explica el error y hazlo bien.

\ex{3} Un huerto cuadrado tiene una superficie de 1225\,m². Lo queremos vallar con 3 vueltas de alambre, dejando una puerta de 5\,m. El alambre se vende en rollos de 50\,m a 18\,€ el rollo. ¿Cuántos rollos hay que comprar? ¿Cuánto costarán?


% =====================================================
\seccio{2}{Números enteros. Valor absoluto}

\ex{1} Ordena de menor a mayor: $\ {-15},\ 8,\ {-3},\ 0,\ {-21},\ 5,\ {-8},\ 12$.

\ex{1} Calcula:\begin{tasks}[after-item-skip=4pt](3)
  \task $\left\lvert -15\right\rvert +\left\lvert 7\right\rvert $
  \task $\left\lvert -4-9\right\rvert $
  \task $\left\lvert -4\right\rvert -\left\lvert 9\right\rvert $
  \task $-\left\lvert -12\right\rvert $
  \task $\left\lvert 3-11\right\rvert \cdot \left\lvert -2\right\rvert $
  \task $\left\lvert -20\right\rvert :\left\lvert -5\right\rvert -\left\lvert -3\right\rvert $
\end{tasks}\recorda{El valor absoluto $\lvert a\rvert $ es la distancia de $a$ al 0: siempre es positivo o cero. Si dentro hay operaciones, primero las hacemos y después quitamos el valor absoluto.}

\ex{2} Escribe todos los números enteros $x$ que cumplen:\begin{tasks}(4)
  \task $\lvert x\rvert =9$
  \task $\lvert x\rvert \le3$
  \task $-4\le x<2$
  \task $\lvert x-2\rvert =5$
\end{tasks}

\ex{3} Calcula paso a paso:\begin{tasks}[after-item-skip=6pt](2)
  \task $\left\lvert 4-9\right\rvert -\left\lvert 2-3\cdot 5\right\rvert +\left\lvert -6\right\rvert \cdot \left(-2\right)$
  \task $\left(3-8\right)\cdot 2+\left\lvert 6-4\cdot 3\right\rvert \cdot \left\lvert -1-2\right\rvert $
  \task $7-\left\lvert -3-\left\lvert 2-9\right\rvert \right\rvert $
  \task $\left\lvert \left(-2\right)\cdot \left(-5\right)-\left\lvert 3-15\right\rvert \right\rvert -\left\lvert -4\right\rvert \cdot \left(1-3\right)$
\end{tasks}

\ex{2} A las 7 de la mañana, en una ciudad de Siberia, el termómetro marcaba $-27$\,°C. A mediodía la temperatura había subido 19 grados y a medianoche había bajado 14 grados respecto del mediodía. a) ¿Qué temperatura hacía a mediodía y a medianoche? b) ¿Qué diferencia hay entre la temperatura más alta y la más baja de estas tres?

\ex{2} Arquímedes nació el año 287 a.\,C. y murió el año 212 a.\,C. ¿Cuántos años vivió?


% =====================================================
\seccio{3}{Operaciones combinadas con enteros}

\ex{1} Calcula:\begin{tasks}(3)
  \task $-12+7-\left(-5\right)$
  \task $\left(-8\right)\cdot \left(-6\right):\left(-4\right)$
  \task $\left(-72\right):\left(-9\right)\cdot \left(-2\right)$
  \task $-15-\left(-3\right)\cdot 4$
  \task $\left(-2\right)\cdot \left(-3\right)\cdot \left(-5\right)$
  \task $-36:\left(-4\right)-9$
\end{tasks}\recorda{Regla de los signos: $(+)\cdot(+)=+$, \ $(-)\cdot(-)=+$, \ $(+)\cdot(-)=-$, \ $(-)\cdot(+)=-$. También para la división.}

\ex{2} Calcula paso a paso:\begin{tasks}[after-item-skip=6pt](2)
  \task $12-4\cdot \left[6-3\cdot \left(5-8\right)\right]$
  \task $-3+24:\left[-2-2\cdot \left(1-6\right)\right]$
  \task $5\cdot \left(-4+9\right)-\left[7-\left(-3\right)\cdot \left(2-6\right)\right]$
  \task $\left(-2\right)\cdot \left[8-\left(3-5\right)\cdot \left(-4\right)\right]-\left(-6\right):\left(-2\right)$
  \task $40:\left[-12+4\cdot \left(7-9\right)\cdot \left(-1\right)\right]$
  \task $-\left[6-\left(2-11\right)\right]+3\cdot \left[-4+\left(5-1\right)\cdot 2\right]$
\end{tasks}

\ex{2} Calcula paso a paso (con potencias y raíces):\begin{tasks}[after-item-skip=6pt](2)
  \task $\left(-3\right)^{3}+\left(-2\right)^{4}-\left(-1\right)^{9}$
  \task $-4^{2}+\left(-4\right)^{2}-\sqrt{36}\cdot \left(-2\right)^{3}$
  \task $\left[\left(-5\right)^{2}-\sqrt{100-36}\cdot 3\right]:\left(-1\right)^{7}$
  \task $2\cdot \left(-3\right)^{2}-\left[\left(-2\right)^{3}+\sqrt{81}\right]^{2}$
  \task $\left(-2\right)^{5}:\left(-2\right)^{3}-3\cdot \left(-1\right)^{20}+\sqrt{4^{2}+3^{2}}$
\end{tasks}\recorda{$(-4)^2=16$, pero $-4^2=-16$: sin paréntesis, el signo no forma parte de la base.}

\ex{2} Una cuenta bancaria tiene un saldo de 420\,€. Se pagan 3 recibos de 85\,€ y 2 de 120\,€ y, después, se ingresan dos nóminas de 150\,€. Escríbelo en una sola expresión y calcula el saldo final. ¿Ha llegado a estar en negativo?

\ex{3} Un ascensor está en la planta 3. Baja 5 plantas, después sube el doble de las que ha bajado y, finalmente, baja 9 plantas. ¿En qué planta está? Si entre planta y planta hay 3\,m, ¿cuántos metros ha recorrido en total?

\ex{3} Las temperaturas mínimas de una semana han sido $-6$, $3$, $-2$, $-8$, $1$, $-4$ y $2$\,°C. Calcula la temperatura media y la diferencia entre la máxima y la mínima.

\ex{3} Halla dos números enteros que multiplicados den $-24$ y sumados den $-2$.


% =====================================================
\seccio{4}{Fracciones. Fracciones equivalentes}

\ex{1} Simplifica hasta obtener la fracción irreducible:\begin{tasks}[after-item-skip=8pt](6)
  \task $\frac{84}{140}$
  \task $\frac{126}{315}$
  \task $\frac{210}{462}$
  \task $\frac{360}{840}$
  \task $\frac{156}{234}$
  \task $\frac{275}{1100}$
\end{tasks}\recorda{Para llegar directamente a la fracción irreducible, dividimos el numerador y el denominador entre su MCD.}

\ex{1} Calcula:\begin{tasks}(1)
  \task Los $\frac{2}{3}$ de los $\frac{3}{5}$ de 450\,€.
  \task Los $\frac{3}{4}$ de la mitad de 96 alumnos.
  \task Los $\frac{5}{6}$ de un tercio de 720 litros.
\end{tasks}

\ex{2} Comprueba si son equivalentes:\begin{tasks}[after-item-skip=6pt](3)
  \task $\frac{21}{35}$ y $\frac{33}{55}$
  \task $\frac{18}{24}$ y $\frac{27}{32}$
  \task $\frac{-14}{49}$ y $\frac{6}{-21}$
\end{tasks}

\ex{2} Halla el valor de $x$ para que las fracciones sean equivalentes:\begin{tasks}[after-item-skip=6pt](3)
  \task $\frac{x}{18}=\frac{35}{42}$
  \task $\frac{12}{x}=-\frac{4}{7}$
  \task $\frac{-15}{25}=\frac{x}{-40}$
\end{tasks}

\ex{2} ¿Qué fracción del día son 16 horas? ¿Y 45 minutos? ¿Y 2 h 40 min? Dalas irreducibles.

\ex{2} Ordena de menor a mayor y sitúalas aproximadamente en la recta: $\ \frac{5}{6},\ -\frac{3}{4},\ \frac{7}{9},\ -\frac{2}{3},\ \frac{13}{18},\ 0$.

\ex{3} Escribe una fracción comprendida entre $\frac{3}{7}$ y $\frac{4}{7}$, y otra entre $\frac{2}{5}$ y $\frac{1}{2}$. Explica cómo lo has hecho.


% =====================================================
\seccio{5}{Operaciones con fracciones}

\ex{1} Calcula y simplifica el resultado:\begin{tasks}[after-item-skip=8pt](4)
  \task $\frac{5}{12}+\frac{7}{18}$
  \task $\frac{11}{15}-\frac{7}{10}$
  \task $3-\frac{11}{4}+\frac{5}{6}$
  \task $-\frac{7}{20}+\frac{5}{8}$
\end{tasks}\recorda{Para sumar o restar, primero reducimos a común denominador (con el mcm de los denominadores). ¡Nunca sumamos los denominadores!}

\ex{1} Calcula y simplifica el resultado:\begin{tasks}[after-item-skip=8pt](4)
  \task $\frac{14}{15}\cdot \frac{25}{21}$
  \task $-\frac{9}{16}\cdot \frac{8}{27}$
  \task $\frac{12}{35}:\frac{18}{25}$
  \task $\left(-\frac{5}{6}\right):\left(-\frac{10}{9}\right)$
\end{tasks}\recorda{$\frac{a}{b}\cdot\frac{c}{d}=\frac{a\cdot c}{b\cdot d}$ \qquad $\frac{a}{b}:\frac{c}{d}=\frac{a\cdot d}{b\cdot c}$. Simplifica antes de multiplicar y los números serán más pequeños.}

\ex{2} Halla la fracción resultante:\begin{tasks}[after-item-skip=6pt](4)
  \task $\frac{2}{3}$ de $\frac{9}{10}$
  \task $la mitad$ de $\frac{4}{7}$
  \task $\frac{3}{5}$ de $\frac{5}{12}$
  \task $\frac{7}{8}$ de $\frac{4}{21}$
\end{tasks}

\ex{2} Calcula (cuidado con los signos):\begin{tasks}[after-item-skip=8pt](2)
  \task $-\frac{5}{6}-\left(-\frac{3}{4}\right)+-\frac{1}{3}$
  \task $\left(-\frac{2}{3}\right)\cdot \frac{9}{4}\cdot \left(-\frac{8}{15}\right)$
  \task $\left(-\frac{4}{9}\right):\left(-\frac{2}{3}\right)\cdot \left(-6\right)$
  \task $-2-\left(-\frac{5}{4}\right)+\frac{7}{6}$
\end{tasks}

\ex{2} Escribe el opuesto y el inverso de cada número:\begin{tasks}(4)
  \task $\frac{3}{8}$
  \task $-\frac{5}{2}$
  \task $-6$
  \task $1+\frac{1}{3}$
\end{tasks}


% =====================================================
\seccio{6}{Operaciones combinadas con fracciones}

\ex{1} Calcula y compara los resultados de a) y b), y de c) y d):\begin{tasks}[after-item-skip=8pt](4)
  \task $\frac{2}{3}+\frac{1}{2}\cdot \frac{4}{9}$
  \task $\left(\frac{2}{3}+\frac{1}{2}\right)\cdot \frac{4}{9}$
  \task $\frac{3}{4}-\frac{1}{4}:\frac{3}{8}$
  \task $\left(\frac{3}{4}-\frac{1}{4}\right):\frac{3}{8}$
\end{tasks}

\ex{2} Calcula paso a paso:\begin{tasks}[after-item-skip=10pt](2)
  \task $\frac{3}{4}-\left[\frac{1}{2}-\left(\frac{2}{5}-\frac{1}{10}\right)\right]$
  \task $2-\left[\frac{1}{3}+\left(\frac{5}{6}-1\right)\right]+\frac{1}{4}$
  \task $1-\left(\frac{2}{3}+\frac{1}{4}-\frac{1}{2}\right):\frac{5}{6}$
  \task $\frac{4}{5}\cdot \left[\frac{1}{2}-\left(\frac{3}{4}-\frac{1}{8}\right)\right]$
  \task $\left[\frac{2}{3}\cdot \left(\frac{1}{4}+\frac{1}{2}\right)\right]:\left[\left(-3\right)\cdot \left(\frac{1}{6}-\frac{1}{2}\right)\right]$
  \task $\left(\frac{5}{6}-\frac{2}{3}-\frac{1}{12}\right):\left(1-\frac{3}{4}\right)$
\end{tasks}

\ex{2} Calcula paso a paso (con potencias):\begin{tasks}[after-item-skip=8pt](3)
  \task $\left(\frac{3}{2}\right)^{2}-\frac{5}{4}\cdot \left(1-\frac{1}{5}\right)$
  \task $\left(-\frac{2}{3}\right)^{3}:\frac{4}{9}+\frac{1}{2}$
  \task $\left[1-\left(\frac{1}{2}\right)^{2}\right]^{2}\cdot \frac{16}{3}$
\end{tasks}\recorda{Para elevar una fracción, elevamos el numerador y el denominador: $\left(-\frac{2}{3}\right)^3=-\frac{2^3}{3^3}=-\frac{8}{27}$.}

\ex{3} Calcula (primero el numerador y el denominador, y después divide):\begin{tasks}[after-item-skip=10pt](4)
  \task $\dfrac{\frac{2}{3}-\frac{1}{4}}{\frac{5}{6}+\frac{1}{3}}$
  \task $\dfrac{1-\frac{3}{5}\cdot \frac{5}{6}}{\left(\frac{1}{2}\right)^{2}+\frac{1}{3}}$
  \task $\dfrac{\left(-2\right)\cdot \left(\frac{1}{4}-\frac{2}{3}\right)}{3\cdot \left(\frac{5}{6}-\frac{1}{2}\right)}$
  \task $\dfrac{\frac{3}{4}:\frac{9}{8}-\frac{1}{3}}{1-\frac{4}{9}}$
\end{tasks}

\ex{3} Calcula paso a paso:\begin{tasks}[after-item-skip=8pt](1)
  \task $\left[\frac{3}{4}-\left(\frac{1}{2}\right)^{2}\right]^{2}:\left(-\frac{1}{2}\right)^{3}+\frac{5}{3}\cdot \frac{6}{5}$
  \task $\frac{1}{2}-\frac{2}{3}\cdot \left[\frac{3}{4}-\left(\frac{1}{2}-2\right)^{2}:3\right]$
\end{tasks}


% =====================================================
\seccio{7}{Problemas con fracciones}

\ex{1} Los $\frac{4}{9}$ de los alumnos de un instituto, que son 352, vienen en autobús. ¿Cuántos alumnos tiene el instituto?

\ex{2} Un móvil cuesta 480\,€. Pagamos $\frac{1}{4}$ de entrada y el resto en 9 cuotas iguales. ¿Cuánto pagamos en cada cuota? ¿Qué fracción del precio es cada cuota?

\ex{2} Se reparten 1800\,€ entre tres hermanos: el primero recibe $\frac{2}{5}$ del total, el segundo $\frac{3}{4}$ de lo que queda y el tercero el resto. ¿Cuánto recibe cada uno? ¿Qué fracción del total recibe el tercero?

\ex{2} En un campo de 2400\,m² se plantan tomates en $\frac{3}{8}$ de la superficie y patatas en $\frac{2}{5}$ del resto. El resto es para pimientos. ¿Cuántos m² ocupa cada cultivo? ¿Qué fracción del campo ocupan los pimientos?

\ex{3} Un grifo llena $\frac{1}{6}$ de un depósito cada hora y otro, $\frac{1}{4}$. Si los abrimos los dos a la vez, ¿qué fracción del depósito se llena en una hora? ¿Cuánto tiempo tardan en llenarlo?

\ex{3} Anna gasta $\frac{2}{7}$ del sueldo en el alquiler y $\frac{1}{3}$ de lo que le queda en comida. Todavía le quedan 600\,€. ¿Cuánto cobra?

\ex{3} En un viaje, el primer día hacemos $\frac{2}{5}$ del recorrido y el segundo día, $\frac{3}{4}$ de lo que falta. Todavía quedan 90\,km. ¿Cuántos kilómetros tiene el viaje?

\ex{3} Un bote lleno de miel pesa 950\,g. Cuando hemos gastado $\frac{3}{5}$ de la miel, pesa 470\,g. ¿Cuánto pesa el bote vacío?

\ex{3} De un depósito lleno, el lunes se gasta $\frac{1}{3}$; el martes, $\frac{1}{4}$ de lo que quedaba, y el miércoles, $\frac{2}{5}$ de lo que quedaba. Al final quedan 360 litros. ¿Cuál es la capacidad del depósito?


% =====================================================
\seccio{8}{Potencias con base entera y racional}

\ex{1} Calcula:\begin{tasks}(6)
  \task $\left(-4\right)^{3}$
  \task $-4^{2}$
  \task $\left(-1\right)^{23}$
  \task $-\left(-2\right)^{4}$
  \task $\left(-10\right)^{5}$
  \task $\left(-9\right)^{0}$
\end{tasks}

\ex{1} Calcula:\begin{tasks}[after-item-skip=8pt](6)
  \task $\left(\frac{3}{4}\right)^{3}$
  \task $\left(-\frac{2}{3}\right)^{4}$
  \task $\left(-\frac{1}{5}\right)^{3}$
  \task $-\left(\frac{3}{2}\right)^{2}$
  \task $\left(-\frac{7}{8}\right)^{0}$
  \task $\left(-\frac{1}{10}\right)^{5}$
\end{tasks}\recorda{$\left(\frac{a}{b}\right)^n=\frac{a^n}{b^n}$. Si la base es negativa, el resultado es positivo con exponente par y negativo con exponente impar.}

\ex{2} Sin hacer el cálculo, di qué signo tiene el resultado:\begin{tasks}(5)
  \task $(-6)^{9}$
  \task $(-11)^{12}$
  \task $-3^{8}$
  \task $\left(-\frac{2}{5}\right)^{11}$
  \task $-(-1)^{100}$
\end{tasks}

\ex{2} Calcula:\begin{tasks}[after-item-skip=8pt](2)
  \task $\left(-3\right)^{3}-\left(-3\right)^{2}-3^{2}$
  \task $\left(-1\right)^{15}-\left(-1\right)^{16}+\left(-2\right)^{3}$
  \task $\left(\frac{2}{3}\right)^{2}-\left(\frac{1}{3}\right)^{2}$
  \task $\left(-\frac{1}{2}\right)^{3}\cdot 8-\left(\frac{3}{4}\right)^{2}:\frac{9}{8}$
\end{tasks}

\ex{3} Dejamos caer una pelota desde 16\,m de altura. Cada vez que bota, sube $\frac{3}{4}$ de la altura anterior. a) ¿A qué altura llega después del cuarto bote? b) ¿Cuántos metros ha recorrido en total cuando toca el suelo por tercera vez?


% =====================================================
\seccio{9}{Propiedades de las potencias}

\ex{1} Expresa como una sola potencia:\begin{tasks}(3)
  \task $5^{3}\cdot 5^{4}$
  \task $2^{9}:2^{4}$
  \task $\left(3^{2}\right)^{5}$
  \task $4^{3}\cdot 5^{3}$
  \task $18^{6}:6^{6}$
  \task $\left(-7\right)^{4}\cdot \left(-7\right)^{3}$
\end{tasks}\recorda{Misma base: $a^m\cdot a^n=a^{m+n}$, \ $a^m:a^n=a^{m-n}$. \ Potencia de potencia: $(a^m)^n=a^{m\cdot n}$. \ Mismo exponente: $a^n\cdot b^n=(a\cdot b)^n$, \ $a^n:b^n=(a:b)^n$.}

\ex{1} Expresa como una sola potencia:\begin{tasks}[after-item-skip=8pt](2)
  \task $\left(\frac{2}{5}\right)^3\cdot\left(\frac{2}{5}\right)^4$
  \task $\left(\frac{3}{4}\right)^8:\left(\frac{3}{4}\right)^5$
  \task $\left[\left(-\frac{1}{2}\right)^3\right]^2$
  \task $\left(\frac{2}{3}\right)^4\cdot\left(\frac{9}{4}\right)^4$
\end{tasks}

\ex{2} Simplifica aplicando las propiedades y calcula:\begin{tasks}[after-item-skip=8pt](3)
  \task $\left(-2\right)^{7}:\left(-2\right)^{4}$
  \task $\left[\left(-1\right)^{3}\right]^{5}\cdot \left(-1\right)^{2}$
  \task $2^{5}\cdot 2^{3}:2^{6}$
  \task $\left(\frac{2}{3}\right)^{5}\cdot \left(\frac{2}{3}\right)^{2}:\left(\frac{2}{3}\right)^{4}$
  \task $\left(-\frac{1}{3}\right)^{6}:\left(-\frac{1}{3}\right)^{4}$
\end{tasks}

\ex{2} ¿Verdadero o falso? Corrige las falsas:\begin{tasks}(3)
  \task $2^3+2^4=2^7$
  \task $3^2\cdot3^3=3^5$
  \task $(2+3)^2=2^2+3^2$
  \task $(3^2)^3=3^5$
  \task $5^0=0$
  \task $6^4:3^4=2^4$
\end{tasks}

\ex{3} Calcula descomponiendo las bases en factores primos:\begin{tasks}[after-item-skip=8pt](3)
  \task $\dfrac{10^4\cdot2^3}{5^3\cdot4^5}$
  \task $9^3\cdot27:3^7$
  \task $\dfrac{12^3\cdot6^2}{18^2\cdot2^5}$
\end{tasks}

\ex{3} Calcula de la manera más rápida: $\ \left(\frac{4}{9}\right)^2\cdot\left(\frac{3}{2}\right)^4$


% =====================================================
\seccio{10}{Potencias de exponente negativo}

\ex{1} Calcula el valor de estas potencias:\begin{tasks}[after-item-skip=8pt](4)
  \task $5^{-2}$
  \task $2^{-5}$
  \task $10^{-3}$
  \task $\left(-3\right)^{-2}$
  \task $\left(-2\right)^{-3}$
  \task $\left(\frac{1}{4}\right)^{-2}$
  \task $\left(\frac{2}{5}\right)^{-3}$
  \task $\left(-\frac{3}{2}\right)^{-1}$
\end{tasks}\recorda{$a^0=1$, \ $a^{-n}=\dfrac{1}{a^n}$ \ y \ $\left(\dfrac{a}{b}\right)^{-n}=\left(\dfrac{b}{a}\right)^{n}$ \ (si $a\neq0$). El exponente negativo no cambia el signo del resultado: \ $2^{-3}=\frac{1}{8}$, no $-8$.}

\ex{2} Ordena de menor a mayor: $\ (-2)^{-1},\ 2^{-2},\ (-2)^{-3},\ 2^{0},\ (-2)^{2},\ 2^{-1}$.

\ex{2} Expresa como una potencia de base 2:\begin{tasks}[after-item-skip=8pt](4)
  \task $8^{-2}$
  \task $\left(\frac{1}{4}\right)^{3}$
  \task $\left(\frac{1}{16}\right)^{-2}$
  \task $\dfrac{4^{-1}\cdot2^{5}}{8^{2}}$
\end{tasks}

\ex{2} Expresa como una sola potencia (con exponentes enteros):\begin{tasks}[after-item-skip=8pt](3)
  \task $7^{-3}\cdot7^{5}$
  \task $3^{4}:3^{-2}$
  \task $(5^{-2})^{3}$
  \task $2^{-4}\cdot5^{-4}$
  \task $\left(\frac{3}{4}\right)^{-2}\cdot\left(\frac{3}{4}\right)^{5}$
  \task $6^{-3}:2^{-3}$
\end{tasks}

\ex{3} Reduce descomponiendo en factores primos y calcula:\begin{tasks}[after-item-skip=8pt](3)
  \task $\dfrac{9^{2}\cdot3^{-5}}{27^{-1}}$
  \task $\dfrac{4^{3}\cdot6^{-2}}{3^{-4}\cdot2^{5}}$
  \task $\dfrac{10^{-2}\cdot25^{2}}{5^{3}\cdot2^{-3}}$
\end{tasks}

\ex{3} Calcula paso a paso:\begin{tasks}[after-item-skip=10pt](2)
  \task $2^{-1}+2^{-2}-2^{-3}$
  \task $\left(\frac{2}{3}\right)^{-2}-\left(\frac{1}{2}\right)^{-1}\cdot \left(-3\right)^{-1}$
  \task $\left[\left(\frac{1}{2}\right)^{-2}-\left(-3\right)^{0}\right]^{-1}$
  \task $\left(\frac{3}{4}-\frac{1}{2}\right)^{-2}:\left(2-\frac{4}{3}\right)^{-1}$
\end{tasks}

\vfill
\begin{center}\color{muted}\small\faCheckCircle[regular]\ \ ¡Ánimo! Comprueba tus resultados con las herramientas interactivas de los apuntes de la unidad.\end{center}

\end{document}
